\section{Spectral projections and block gradings}
\label{spec:section}

The shift operator $T_a$ distinguishes $r+1$ spectral blocks at generic
fiber parameter, whereas the localization formula distinguishes the two fixed
manifolds $B$ and $X$.  We construct operators on the individual spectral
blocks and compare their sum near $y=0$ with the fixed-manifold blocks.
The construction must also remove differentiation in the equivariant
parameter: only after that removal can the resulting loop operators be
quantized over the coefficient ring.  All assertions in this section are
classical statements about operators.

\subsection{Functional calculus for differential series}

Let $V$ be the finite-dimensional vector space underlying a chosen
equivariant basis of $H_T^*(W)$.  On a small open set $U$ in the
$(y,\lambda)$-plane, write $\mathcal O_U$ for holomorphic functions and
\[
 \mathscr D_\lambda(U)
 =\left\{\sum_{j=0}^{J}a_j(y,\lambda)\partial_\lambda^j:
                 J<\infty,\ a_j\in\mathcal O_U\right\}.
\]
The multiplication is composition, so
$\partial_\lambda a=a\partial_\lambda+\partial_\lambda(a)$.
There is no differentiation in $y$.  Denote by $\Lambda_B$ the base
Novikov completion already fixed, and by $\Lambda_{B,+}$ its positive
degree ideal.  We use the completed algebra
\begin{equation}
 \mathscr A_U
 =\left\{\sum_{\beta,m\geq0}Q^\beta z^m A_{\beta,m}:
      A_{\beta,m}\in\End(V)\otimes\mathscr D_\lambda(U)\right\}.
 \label{spec:algebra}
\end{equation}
Here and below the notation $\beta\geq0$ means that $\beta$ ranges over
the effective base labels, including zero.  The completion is taken with
respect to $z$ and ample base degree.  In particular each
$A_{\beta,m}$ has finite differential order, without a bound uniform in
$\beta,m$.  Products at a fixed coefficient are finite by Novikov
finiteness.  The element $z$ is central in this algebra.  We may equally
work with germs, and then all holomorphic assertions are coefficientwise.
When a loop derivative occurs, we adjoin $z\partial_z$ with
\[
 [z\partial_z,Q^\beta z^m A_{\beta,m}]
                =mQ^\beta z^m A_{\beta,m};
\]
all expressions involving this extra derivative have finite order in it.

Suppose $T\in\mathscr A_U$ has order-zero reduction
\[
 T^{(0)}=T\bmod(z,\Lambda_{B,+})\in
                \End(V)\otimes\mathcal O_U.
\]
After shrinking $U$, choose contours in the complex $\xi$-plane that
avoid the spectrum of $T^{(0)}$ and enclose specified groups of its
eigenvalues.  The contours are held fixed while $y$ and $\lambda$ vary
in $U$.  Set $R_0(\xi)=(\xi-T^{(0)})^{-1}$ and
$K=T-T^{(0)}$.  The formal resolvent is
\begin{equation}
 R_T(\xi)=(\xi-T)^{-1}
    =\sum_{n\geq0}\bigl(R_0(\xi)K\bigr)^nR_0(\xi).
 \label{spec:resolvent}
\end{equation}
Each factor $K$ increases the joint $z$/base filtration.  Consequently
each coefficient in this expression is a finite sum of differential
operators, and is meromorphic in $\xi$, with poles only at the
eigenvalues of $T^{(0)}$.

\begin{lemma}[Coefficientwise holomorphic calculus]
\label{spec:calculus}
For $T$ as above and a scalar holomorphic function $f$ on a neighborhood of
$\Spec(T^{(0)})$, held fixed independently of the parameters, define
\begin{equation}
 f(T)=\frac{1}{2\pi\mathrm i}\int_\Gamma
                  f(\xi)R_T(\xi)\,d\xi,
 \label{spec:functional-calculus}
\end{equation}
where $\Gamma$ encloses the whole spectrum within that neighborhood.
The neighborhood may be disconnected.  This construction is unital,
multiplicative, and compatible with holomorphic composition.  In
particular, the functions equal to $1$ on one spectral group and $0$
on the others give mutually orthogonal idempotents whose sum is the
identity.  All these operators commute with $T$.

If, in addition, $[T,\lambda]=zT$, then
\begin{equation}
 [f(T),\lambda]=z\bigl(\xi f'(\xi)\bigr)(T).
 \label{spec:functional-commutator}
\end{equation}
If a differential operator $G$ commutes with $T$, acts continuously on
the coefficientwise Laurent extension of \eqref{spec:algebra}, and
commutes with the scalar spectral variable $\xi$, then
$[G,f(T)]=0$.
\end{lemma}

\begin{proof}
Both inverse identities for \eqref{spec:resolvent} follow from the
geometric series.  Since $\xi$ and a second spectral variable $\eta$
are central, the resolvent identity is
\[
 R_T(\xi)R_T(\eta)
       =\frac{R_T(\xi)-R_T(\eta)}{\eta-\xi}.
\]
Integrating on nested contours and applying the scalar Cauchy formula
proves multiplicativity.  It also proves $1(T)=\id$, either directly
or coefficientwise from \eqref{spec:resolvent}.  This gives the
projector assertions.  Commutation with $T$ follows from the inverse
identity.

For completeness, multiplicativity also gives the composition rule.
For $\eta$ outside the spectrum of $f(T^{(0)})$, apply the calculus to
$(\eta-f(\xi))^{-1}$.  Its product with $\eta-f(T)$ is the identity,
so it is the resolvent of $f(T)$.  Integrating this identity against
$g(\eta)$ and using the scalar Cauchy formula yields
$g(f(T))=(g\circ f)(T)$ whenever the functions are defined on the
indicated neighborhoods.  All contour operations here are performed
on a fixed coefficient, where there are only finitely many
differential compositions.

The commutator with $\lambda$ follows directly from the inverse rule:
\begin{align*}
 [R_T(\xi),\lambda]
   &=R_T(\xi)[T,\lambda]R_T(\xi)\\
   &=z\bigl(\xi R_T(\xi)^2-R_T(\xi)\bigr)
     =-z\partial_\xi\bigl(\xi R_T(\xi)\bigr).
\end{align*}
Integration by parts proves \eqref{spec:functional-commutator}.
Finally, $[G,T]=[G,\xi]=0$ implies $[G,R_T(\xi)]=0$ by the
inverse identity.  Integration proves the last assertion.  Keeping
the contours fixed locally justifies differentiating the coefficient
germs under the integral.
\end{proof}

We will compare this calculus in frames related by matrices with
negative powers of $z$.  The comparison requires a larger algebra,
but does not require convergence in the loop variable.  Let $\mathcal
O$ be a ring of holomorphic parameter germs stable under
$\partial_\lambda$, and let $\mathcal M$ be either the base Novikov
monoid or the full monoid of labels $(\beta,l)$, $l\geq0$.  Write
\begin{equation}
 \mathscr A^{\mathrm{Laur}}_{\mathcal M}(\mathcal O)
 =\left\{\sum_{d\in\mathcal M}q^d
       \sum_{m\geq m(d)}z^m A_{d,m}:
       A_{d,m}\in\End(V)\otimes\mathscr D_\lambda(\mathcal O),
       \ m(d)\in\Z\right\}.
 \label{spec:laurent-algebra}
\end{equation}
The notation $q^d$ means $Q^\beta$ or $Q^\beta y^l$, respectively.
The Novikov support condition is imposed as before.  The lower
$z$-bound can depend on $d$.  Multiplication is well-defined: a fixed
label has finitely many decompositions, and for each decomposition
the two Laurent lower bounds leave only finitely many summands at a
fixed power of $z$.

\begin{lemma}[Change of frame for resolvents]
\label{spec:gauge}
Let $T$ and $T'$ have coefficientwise resolvents on a common contour,
constructed as in Lemma~\ref{spec:calculus}.  Suppose these resolvents
belong to the same algebra \eqref{spec:laurent-algebra}.  If $M$ and
$M^{-1}$ belong to that algebra, are independent of $\xi$, and
$T'=MTM^{-1}$, then
\[
 R_{T'}(\xi)=M R_T(\xi)M^{-1},\qquad
 f(T')=M f(T)M^{-1}
\]
for every function represented by that contour calculus.
\end{lemma}

\begin{proof}
The expression $M R_T(\xi)M^{-1}$ is both a left and a right inverse
of $\xi-T'$ in the common algebra.  It therefore equals its other
inverse $R_{T'}(\xi)$.  Integrating the equality coefficientwise
gives the second assertion.  The Laurent lower bounds ensure that
each product used in the integration is finite at a fixed
coefficient.
\end{proof}

Two forms of $M$ will occur.  The first is a matrix equal to the
identity in degree zero whose coefficient at each positive Novikov
label is the Laurent expansion at $z=0$ of a rational function of $z$.
Its Novikov inverse has the same property.  The second is a matrix
of the form $\exp(N/z)M_+(z)$, where $N$ is a nilpotent cup-product
operator and $M_+(z)$ and its inverse are regular at $z=0$.
The exponential and its inverse are then Laurent polynomials in
$z^{-1}$.  Thus both forms satisfy the Laurent requirement of
Lemma~\ref{spec:gauge}.

\subsection{Block operators without equivariant differentiation}

We now apply the calculus to $T=T_a$.  By
Lemma~\ref{shift:polynomiality}, its coefficients are polynomial in
$y,\lambda,z$ at fixed base degree before the shift
$\exp(z\partial_\lambda)$ is expanded.  In particular it belongs
to \eqref{spec:algebra}.  By \eqref{shift:spectrum}, the spectrum of
its reduction modulo $z$ and positive base degree consists of
\begin{equation}
 a_i=\lambda+h_i,\qquad h_i^r(h_i+\lambda)=y.
 \label{spec:eigenvalues}
\end{equation}
Take $\lambda\neq0$, $y\neq0$, and exclude the values where these
roots coincide.  On a small open set in this locus, choose their
labels and a logarithm near each $a_i$.  A label refers to a whole
generalized eigenspace: no diagonalizability within that space is
assumed.  Define
\begin{equation}
 P_i=\frac{1}{2\pi\mathrm i}\int_{\gamma_i}R_{T_a}(\xi)\,d\xi,
 \qquad
 \mathscr L_i=\frac{1}{2\pi\mathrm i}\int_{\gamma_i}
                  \log\xi\,R_{T_a}(\xi)\,d\xi.
 \label{spec:projectors}
\end{equation}
Here $\gamma_i$ encloses the one spectral value $a_i$ with its full
multiplicity.  For a group of these values we use the sum of the
contours; its logarithm is specified on each component.

\begin{proposition}[Removal of coefficient differentiation]
\label{spec:commutators}
The operators in \eqref{spec:projectors} satisfy
\begin{equation}
 [P_i,\lambda]=0,\qquad
 [\mathscr L_i,\lambda]=zP_i,\qquad
 [G_a,P_i]=[G_a,\mathscr L_i]=0.
 \label{spec:block-commutators}
\end{equation}
In particular $P_i$ and
$M_i=\mathscr L_i-zP_i\partial_\lambda$ are matrix series,
without coefficient differentiation.  The operator
\begin{equation}
 D_i=P_iG_a-\frac{\lambda}{z}\mathscr L_i
 \label{spec:block-grading}
\end{equation}
differentiates only in $z$ and satisfies
$P_iD_i=D_iP_i=D_i$.
\end{proposition}

\begin{proof}
The shift form gives $[T_a,\lambda]=zT_a$.
Apply \eqref{spec:functional-commutator} to the locally constant
branch indicator and to its supported logarithm.  These give the
first two commutators.  The relation
$[G_a,T_a]=0$ from \eqref{shift:commutation} gives the last two.
The operator $G_a$ has only one negative loop power and first-order
parameter derivatives, so its action and the commutators used here
are defined in \eqref{spec:laurent-algebra}.

For a finite-order differential operator $D$ in $\lambda$,
$[D,\lambda]=0$ implies that $D$ has order zero: the commutator of
$a_J\partial_\lambda^J$ with $\lambda$ has leading term
$J a_J\partial_\lambda^{J-1}$.  Apply this observation at each
coefficient of $P_i$ and $\mathscr L_i-zP_i\partial_\lambda$.
Writing out $G_a$ now gives
\begin{equation}
 D_i=P_i z\partial_z
       +P_i(\tfrac12+\mu)
       +z^{-1}\bigl(P_i(c\star)-\lambda M_i\bigr).
 \label{spec:grading-matrix-form}
\end{equation}
There are no remaining $\lambda$ derivatives, and neither $G_a$ nor
the functional calculus introduced derivatives in $y$ or $Q$.
The support assertions follow from $[G_a,P_i]=0$,
$P_i\mathscr L_i=\mathscr L_iP_i=\mathscr L_i$, and
$[P_i,\lambda]=0$.
\end{proof}

The subtraction in \eqref{spec:block-grading} has accomplished the
first goal: it gives loop operators over the coefficient ring.
We record their precise bounds before making the comparison at
$y=0$.

\begin{proposition}[Modes and their coefficient bounds]
\label{spec:modes}
Define
\begin{equation}
 A_{-1,i}=z^{-1}P_i,\qquad
 A_{k,i}=z^{-1}(zD_i)^{k+1}\quad(k\geq0).
 \label{spec:mode-definition}
\end{equation}
Each $A_{k,i}$ is a series in powers $z^m$ with $m\geq-1$, whose
coefficients are matrix differential operators in $z\partial_z$
of order at most $k+1$.  Its coefficients contain no derivatives
in $y,\lambda$ or $Q$.  Each fixed base and loop coefficient
continues holomorphically along paths in the generic locus of
\eqref{spec:eigenvalues}, with the spectral labels and logarithms
continued along the path.

For any scalar $b$ independent of $z$, replacing $D_i$ by
$D_i+bz^{-1}P_i$ gives
\begin{equation}
 A_{k,i}\longmapsto
   \sum_{j=0}^{k+1}\binom{k+1}{j}b^{k+1-j}A_{j-1,i}.
 \label{spec:binomial}
\end{equation}
For distinct branches $i\neq j$, one has
$(zD_i)(zD_j)=0$.  Thus, if $P=\sum_iP_i$ and
$D=\sum_iD_i$ for a group of branches, then the modes formed from
$P,D$ are the sums of their individual modes.
\end{proposition}

\begin{proof}
Equation~\eqref{spec:grading-matrix-form} says that $zD_i$ has
nonnegative powers of $z$ and Euler differential order at most
one.  Since $z\partial_z$ preserves each loop power, products do
not create negative powers.  This proves the asserted bounds.

At a fixed base and loop coefficient, the resolvent construction
uses finitely many matrix inversions, parameter derivatives, and
contour residues.  Its only spectral denominators come from the
eigenvalues in \eqref{spec:eigenvalues}.  The input coefficients
are polynomial, so these operations continue along any path on
which the eigenvalues remain distinct and nonzero, with the
logarithms continued.  This is a statement about individual
coefficients; it asserts no convergence of the $z$ or Novikov
series.

Since $D_i$ differentiates only $z$, it commutes with $b$.
Furthermore $P_i$ commutes with $z$ and is a two-sided identity
for $D_i$.  The ordinary binomial formula in this supported
algebra gives \eqref{spec:binomial}, including its $j=0$ term
$b^{k+1}z^{-1}P_i$.  Finally,
\[
 (zD_i)(zD_j)
   =zD_iP_i zP_jD_j=0
 \qquad(i\neq j),
\]
which proves the assertion about sums of branches.
\end{proof}

In particular, changing a logarithm by $2\pi\mathrm i n$ changes
$D_i$ by $-2\pi\mathrm i n\lambda z^{-1}P_i$.  All the modes for
one logarithm therefore span the same collection as those for any
other logarithm.  Infinitesimal symplecticity will follow from the
fixed-manifold calculation and continuation; it is not needed for
the constructions above.

\subsection{The two fixed-manifold blocks at \texorpdfstring{$y=0$}{y=0}}

Fix a germ of $\lambda\neq0$.  At $y=0$, one eigenvalue
$a_B=\lambda+h_B$ tends to zero, while the other $r$ tend to
$\lambda$.  We call the first the $B$ branch and the second group
the $X$ cluster.  Choose disjoint small contours about $0$ and
$\lambda$ that enclose these groups for $y$ sufficiently small.
The projector calculus applies through $y=0$ on both contours.
The logarithm applies through $y=0$ on the $X$ contour, since that
contour bounds a neighborhood not containing zero.  Denote the
resulting operators by $P_B,P_X,\mathscr L_X$.

For this comparison express both frames in the fixed-cohomology
coordinates furnished by restriction to $B\sqcup X$.  The change
from the polynomial equivariant basis depends only on $\lambda$
and is invertible at $\lambda\neq0$; it preserves the completed
algebras and all assertions about holomorphic dependence on $y$.
Let $\Pi_B,\Pi_X$ be the constant projections onto these two
summands.  In these coordinates write
\[
 K_F(z,\lambda)=
 \prod_{x}\begin{cases}
      x+\lambda,&F=X,\\
      (x-\lambda-z)^{-1},&F=B,
 \end{cases}
 \qquad E_z=\exp(z\partial_\lambda),
\]
where the product is over the ordinary Chern roots of $N_F$.
Thus \eqref{shift:fixed-formula} reads
$T_d|_F=y^{j_F}K_F E_z$.

\begin{proposition}[Comparison at zero fiber degree]
\label{spec:zero}
The projectors $P_B,P_X$ and the cluster logarithm
$\mathscr L_X$ have holomorphic coefficient germs through $y=0$.
Their Taylor expansions at $y=0$ satisfy
\begin{equation}
 P_F=S_W\Pi_FS_W^{-1}\quad(F=B,X),\qquad
 \mathscr L_X=S_W\mathscr L_{d,X}S_W^{-1},
 \label{spec:zero-conjugation}
\end{equation}
where $\mathscr L_{d,X}$ is the logarithm of $K_XE_z$ on the $X$
summand, extended by zero on $B$.

Moreover, $T_aP_B$ is divisible by $y$ coefficientwise in $z$ and
base degree.  Set
\begin{equation}
 \widetilde T_B=y^{-1}T_aP_B.
 \label{spec:divided-shift}
\end{equation}
It has holomorphic coefficient germs through $y=0$.  Its
nonzero spectral block at $y=z=Q^{>0}=0$ has eigenvalue
$(-\lambda)^{-r}$, with possible nilpotent part, and its
complementary block is zero.  If $\mathscr K_B$ is its supported
logarithm on the nonzero block, then, on a punctured sector in
$y$, one can choose the $B$-branch logarithm as
\begin{equation}
 \mathscr L_B=(\log y)P_B+\mathscr K_B.
 \label{spec:divided-log}
\end{equation}
The coefficients of $\mathscr K_B$ are holomorphic through $y=0$.
Its Taylor expansion, and hence \eqref{spec:divided-log}, obeys
the same conjugation by $S_W$ as in
\eqref{spec:zero-conjugation}, with the corresponding fixed-block
operators in the descendant frame.
\end{proposition}

\begin{proof}
The reduction $T_a^{(0)}$ is a holomorphic matrix through $y=0$.
Its spectrum there consists of the two separated values
$0,\lambda$, including all their multiplicities and nilpotent
parts.  The Neumann construction with the two fixed contours
therefore gives the asserted holomorphic extensions of
$P_B,P_X,\mathscr L_X$.

To compare frames, first expand these coefficient germs in $y$.
Taylor expansion is a homomorphism commuting with
$\partial_\lambda$; it sends the ancestor resolvent, whose loop
powers are nonnegative, to a Laurent series in the full Novikov labels and
preserves both of its inverse identities.
Use \eqref{spec:laurent-algebra} with full labels
$d=(\beta,l)$, $q^d=Q^\beta y^l$, and coefficients holomorphic
in the chosen $\lambda$ germ.  The rationality assertion in
Proposition~\ref{loc:factorization} puts the expansion of $S_W$ at
$z=0$ in this algebra: a rational function has a finite-order
pole at zero.  Since its zero-curve coefficient is the identity,
its inverse belongs to the same algebra.  The two shift
operators, expanded using \eqref{shift:fixed-formula}, also
belong to it.  Their resolvents on the two contours do as well.
For the descendant resolvent one can use the joint filtration
by $z$ and full Novikov degree: its initial $B$ block is zero,
and its initial $X$ block is multiplication by $\lambda+p$,
whose nilpotent part is the ordinary class $p|_X$.

Equation~\eqref{shift:calibration} and
Lemma~\ref{spec:gauge} now identify the actual resolvents in
the two frames.  In the descendant frame the contour about
zero selects exactly the whole $B$ summand, and the contour
about $\lambda$ selects exactly the whole $X$ summand.
Indeed the reduced spectra have this property; within each
summand the contour function is identically $1$ or $0$, so
Lemma~\ref{spec:calculus} gives respectively the identity or
zero also for the formal perturbation.  Integration proves
\eqref{spec:zero-conjugation}, including the logarithm on $X$.

On the $B$ summand, $T_d\Pi_B=yK_BE_z\Pi_B$.  The parameter
$y$ is central throughout the algebra.  Hence
\[
 T_aP_B
   =yS_WK_BE_z\Pi_BS_W^{-1}
\]
in the Laurent algebra completed in full Novikov degree.  Neither $S_W$ nor its
inverse has a negative $y$ power.  The constant Taylor
coefficient in $y$ of every base and loop coefficient on the
left therefore vanishes.  Those coefficients were already
holomorphic through $y=0$, so division by $y$ gives
holomorphic coefficients and
\begin{equation}
 \widetilde T_B=S_WK_BE_z\Pi_BS_W^{-1}
 \label{spec:divided-conjugation}
\end{equation}
after Taylor expansion.

At $y=z=Q^{>0}=0$, the calibration is the identity and the
nonzero block of \eqref{spec:divided-conjugation} is
$\prod_x(x-\lambda)^{-1}$.  The ordinary positive-degree
classes are nilpotent, so its sole spectral value is
$(-\lambda)^{-r}\neq0$.  Choose a contour enclosing this
value and avoiding zero, and choose a logarithm there.
Lemma~\ref{spec:calculus} defines $\mathscr K_B$ with
holomorphic coefficients through $y=0$.  The projector for
this nonzero block is $P_B$, by
\eqref{spec:divided-conjugation} and
Lemma~\ref{spec:gauge}.  The same lemma identifies
$\mathscr K_B$ with the conjugate of the supported logarithm
of $K_BE_z\Pi_B$.

Finally restrict to a punctured sector and choose $\log y$.
In the algebra supported on $P_B$ one has
$T_a=y\widetilde T_B$.  Rescaling the spectral contour by
the central scalar $y$ and choosing
$\log(y\xi)=\log y+\log\xi$ shows that its logarithm is
exactly \eqref{spec:divided-log}.  This is a choice of the
original contour logarithm in \eqref{spec:projectors}.
The conjugation assertion follows because multiplication
by $\log y$ commutes with the entire algebra.
\end{proof}

We have thus constructed the same block operators in two useful
forms.  At generic $y$, their individual coefficients continue
with the roots in \eqref{spec:eigenvalues}.  Near $y=0$, they
are conjugates of operators on the two fixed manifolds.
The coefficients of the $B$-branch modes are finite polynomials
in $\log y$ with holomorphic coefficients.  Indeed
$D_B=D_B^{\mathrm{reg}}-\lambda(\log y)z^{-1}P_B$, where
$D_B^{\mathrm{reg}}=P_BG_a-(\lambda/z)\mathscr K_B$
has holomorphic coefficients; apply \eqref{spec:binomial} with
$b=-\lambda\log y$.  For $A_{k,B}$ the logarithmic degree
is therefore at most $k+1$.  The $X$-cluster modes are
holomorphic through zero by Proposition~\ref{spec:zero}.
These statements require neither convergence in $z$ nor an
analytic interpretation of the Novikov series.
